\chapter{Three jobs with one name}
\label{ch:three-jobs}

\section*{What this chapter is for}

Three different jobs are advertised under titles that contain the letters
\emph{AI}, and they want different people. Preparing for the wrong one of the
three is the most expensive mistake available before you have spoken to
anybody, because it is invisible: you study hard, you study well, and you study
the wrong surface.

This chapter draws the line between them, says which signals in a posting are
reliable and which are not, and gives you a way to classify a role in about ten
minutes.

\section{The line that actually separates them}

The cleanest statement of the boundary is about ownership of weights.
\reported{Machine-learning engineers own model weights; AI engineers use
third-party models through an API.}{field-guide} The same split is drawn from
the other end in the standard text on the subject: \reported{AI engineering
starts from an available model and prototypes early, and only considers
training a custom one much later, where classical machine learning starts from
data annotation, feature engineering and training.}{aie-book}

That single distinction predicts most of the rest.

\begin{kittable}{@{}lXX@{}}
 & \textbf{AI Engineer} & \textbf{ML Engineer / Research} \\
\midrule
Starts from & a model that exists & data that exists \\
Owns & prompts, retrieval, tools, evals, budgets & weights, training pipelines, features \\
First artefact & a working path end to end & a dataset and a baseline \\
Fails by & the system being wrong in production & the model being wrong on held-out data \\
Is asked about & how you build, evaluate and operate it & how the method works \\
\end{kittable}

The corpus behind the first row of that table summarises its own finding in one
sentence: \reported{interviewers care more about how you build, evaluate and
operate AI systems than about transformer internals.}{field-guide}

\section{What each job owns}

Read these as lists of things you will be held responsible for, not as topics.

\textbf{An AI Engineer owns} prompt design and, more to the point, prompt
\emph{versioning}; the evaluation harness and the datasets it runs on;
retrieval quality and the debugging of it when it is poor; the schema of the
tools a model is allowed to call; production monitoring, logging and the
feedback loop back into the evals; and a cost and latency budget per
request.\source{field-guide}

\textbf{A machine-learning engineer or researcher owns} model weights, the
training and fine-tuning pipeline, feature engineering, dataset curation for
training, and architecture work.\source{field-guide}

There is a third job, thinly advertised and worth naming: the platform role
that builds the gateway, the tracing, the spend controls and the evaluation
infrastructure that the first job consumes. It is scored much more like a
conventional distributed-systems role with one unfamiliar dependency.
\judgement{If your background is platform or infrastructure rather than product
engineering, this is frequently the role you are strongest for and the one you
are least likely to have applied to, because its postings tend to use the
vocabulary of the first job.}

\section{The posting is not the job}

The gap between what postings advertise and what the work turns out to be is
reported bluntly. One hiring manager quoted in the same corpus describes the
role as \reported{eighty per cent prompt engineering, glue code and
monitoring}{field-guide-reality}, against postings that lead on retrieval
architecture and embeddings. The same source records that the daily work is
writing code that \emph{generates} prompts dynamically rather than hand-writing
them, and building systems that run against several providers with fallbacks
between them.

\begin{kitnote}
This is not a complaint about dishonest postings. A posting is written to
attract, and the parts of a job that attract are not the parts that fill the
day. It matters here for one reason only: it tells you which of your existing
experience is relevant. Glue code, monitoring and multi-provider fallback are
things a senior backend engineer has done for a decade under other names.
\end{kitnote}

\section{Classifying a role in ten minutes}

The title is the least reliable signal on the page. Use the verbs.

\begin{kitmethod}
\textbf{1. Find the verbs in the responsibilities section.} \emph{Train},
\emph{fine-tune}, \emph{curate}, \emph{label}, \emph{benchmark against a test
set} point at the second job. \emph{Integrate}, \emph{evaluate}, \emph{ship},
\emph{instrument}, \emph{reduce latency}, \emph{control cost} point at the
first.

\textbf{2. Find out what they would have you touch on day one.} A role that
starts at a dataset is not the same role as one that starts at an endpoint.

\textbf{3. Read the stack.} A posting listing a vector store, an orchestration
library and an observability product is describing the first job whatever the
title says. One listing a training framework and an experiment tracker is
describing the second.

\textbf{4. Check whether evaluation is mentioned at all.} Its presence is a
strong signal the team has shipped something and been hurt by it. Its absence
is not proof of the opposite, but it changes what you should ask in the screen.
\end{kitmethod}

\begin{drill}
Take three current postings that you would plausibly apply to. For each, write
one sentence naming which of the three jobs it is, and one sentence naming the
evidence in the posting that decides it. If you cannot find the evidence, that
is the finding: the posting is ambiguous, and the ambiguity is a question for
the screening call rather than something to resolve by guessing.
\end{drill}

\section*{What this chapter does not claim}

It does not claim the three jobs are cleanly separated in practice. Small teams
collapse all three into one person, and plenty of people with \emph{ML
Engineer} on a badge spend their week on retrieval quality. The boundary is a
tool for deciding what to study and what to ask, not a description of how
anybody's organisation is actually arranged.

It also makes no claim about which job pays better, is more secure, or is more
interesting. Those are questions about a specific company and a specific team,
and Chapter~\ref{ch:company} is about answering them for the one in front of
you rather than in general.
